\subsection{What the natural exterior-volume powers actually detect}
\label{sec:exterior-volume-fermat}

\paragraph{P-F47 (an explicit coefficient map to the Fermat degree monoid).}
Let $q,n\ge1$, put $k=3q$, $d=nq$, and let
$E=\mathbb C^k$ with ordered basis $e_1,\ldots,e_k$.
Write $W=\bigwedge^3E$ and $e_J=e_i\wedge e_j\wedge e_l$
for a triple $J=\{i<j<l\}$. Define
\[
 \theta_k=\sum_{\mathcal M}\epsilon(\mathcal M)
                    e_{J_1}\wedge\cdots\wedge e_{J_q}
       \ \in\ \bigwedge^qW,                         \tag{P-F47a}
\]
where the sum is over the partitions $\mathcal M$ of $[k]$
into triples, their blocks $J_1,\ldots,J_q$ are listed
lexicographically, and $\epsilon(\mathcal M)$ is the sign
of their concatenated increasing lists as a permutation
of $[k]$. Then $\theta_k$ spans a copy of $\det E$, and
\[
 \theta_k^{\wedge n}\ne0
 \quad\Longleftrightarrow\quad
 \exists I\in\mathbb C[\operatorname{Sym}^3\mathbb C^n]^{SL_n}_d
             \text{ with }I(x_1^3+\cdots+x_n^3)\ne0.
                                                        \tag{P-F47b}
\]
More precisely, every coefficient of the exterior power
is, up to an explicitly fixed sign, the evaluation of a
determinant-contraction invariant on the Fermat cubic.
Conversely these coefficients span all evaluations of
degree-$d$ invariants on that cubic.

This is a statement about nonvanishing at one specified
form. It is not an equivalence between this particular
exterior power and nonvanishing of the whole degree-$d$
invariant space. In particular, failure of (P-F47b) does
not imply $f_n(gq)=0$, where $g=\gcd(n,3)$.

\paragraph{The canonical volume inclusion.}
The map $\det E\longrightarrow\bigwedge^q(\bigwedge^3E)$
giving (P-F47a) is explicitly
\[
 v_1\wedge\cdots\wedge v_k\longmapsto
 \frac{1}{q!(3!)^q}\sum_{\sigma\in\mathfrak S_k}
 \operatorname{sgn}(\sigma)
 \bigwedge_{a=0}^{q-1}
 (v_{\sigma(3a+1)}\wedge v_{\sigma(3a+2)}
                          \wedge v_{\sigma(3a+3)}).
                                                        \tag{P-F47c}
\]
It is alternating and multilinear in the $v_i$, so is
well-defined and $GL(E)$-equivariant. Permuting two
blocks contributes a sign $(-1)^9$ in their flattened
lists and a sign $-1$ in the outer wedge; these cancel.
Permuting entries of a triple similarly contributes
two equal signs. Thus each term of (P-F47a) occurs
exactly $q!(3!)^q$ times before division. Its coefficients
are $\pm1$, so the map is nonzero. This proves the
claimed determinant character without choosing a
metric or an unproved splitting of a Schur functor.

\paragraph{Determinant contractions and their values.}
Consider a tableau $T$ with $n$ rows and $k$ columns,
whose entries are labels $1,\ldots,d$, each occurring
three times, with no repeated label in one column.
For each label $a$, let $J_a$ be its three columns.
For a symmetric cubic tensor $F=(F_{ijk})$, define
\[
 I_T(F)=\sum_{\sigma_1,\ldots,\sigma_k\in\mathfrak S_n}
       \left(\prod_{j=1}^k\operatorname{sgn}(\sigma_j)\right)
       \prod_{a=1}^d
       F_{\sigma_{j_1}(r_1),\sigma_{j_2}(r_2),
                                      \sigma_{j_3}(r_3)},
                                                        \tag{P-F47d}
\]
where $(r_b,j_b)$ are the three positions of label $a$.
The order of those positions does not matter because
$F$ is symmetric. Contracting the $n$ indices in each
column with the alternating determinant tensor proves
\[
                    I_T(hF)=(\det h)^k I_T(F).
                                                        \tag{P-F47e}
\]
Thus it is a degree-$d$ polynomial $SL_n$-invariant.
An exact primary construction and transformation proof
are B\"urgisser--Ikenmeyer, \emph{Fundamental invariants
of orbit closures}, equation~(3.7) and Theorem~3.11,
printed pp.~11--12. Their Remark~3.13, p.~13, states
that the semistandard tableaux of this shape and content
span the entire invariant space. Equivalently, one can
polarize a degree-$d$ invariant, use the determinant
contractions spanning the $SL_n$-invariants in the
$3d=nk$ tensor positions, and then symmetrize within
each cubic and among the $d$ cubic factors. A repeated
label in a column gives zero after this symmetrization.

In the tensor convention of (P-F47d), the Fermat cubic
$F_n=\sum_i x_i^3$ has $F_{iii}=1$ and all its other
tensor entries zero. Its nonzero summands are exactly
the colorings $c:\{1,\ldots,d\}\to[n]$ such that in
each column of $T$ every color occurs once. A coloring
determines its $k$ permutations uniquely. Consequently
\[
 I_T(F_n)=\sum_c\prod_{j=1}^k
                  \operatorname{sgn}(c\text{ in column }j).
                                                        \tag{P-F47f}
\]
If $J_a=J_b$ for two distinct labels, exchange their
colors. The colors are different in each of their
three common columns. This is a fixed-point-free
involution on the admissible colorings, and it changes
the product of column signs by $(-1)^3=-1$.
Therefore $I_T(F_n)=0$ whenever two triples coincide.

\paragraph{The exact coefficient sign.}
It remains to consider a set $H=\{J_1,\ldots,J_d\}$
of distinct triples, listed in a fixed order, with
each vertex in exactly $n$ triples. Choose the tableau
by placing the incident labels in any fixed order in
each column. A coloring in (P-F47f) is precisely an
ordered factorization of $H$ into $n$ perfect matchings
of $[k]$. Expansion of $\theta_k^{\wedge n}$ shows that
its coefficient at $e_{J_1}\wedge\cdots\wedge e_{J_d}$
is the signed sum of these same ordered factorizations.
There is no omitted factorial: each matching occurs
once in (P-F47a), and the $n$ factors are ordered.

For an explicit comparison of the signs, distinguish
the $3d=nk$ incidences $(a,j)$ with $j\in J_a$.
Let $A$ be their order by label $a$, then by increasing
vertex $j$ within that triple. Let $B$ be their order
by vertex $j$, then by the chosen tableau row. Write
$\varepsilon_T$ for the sign of this fixed reordering
$A\to B$. Write $\tau_{n,k}$ for the sign of the fixed
transpose from the grid order $(j,i)$ to $(i,j)$.
For a coloring $c$, the sign of reordering $A$ to
the order by color and then vertex can be computed
in two ways. Grouping triples by color has the same
sign as grouping the outer wedge factors, because
interchanging two triples interchanges nine incidences.
Ordering the vertices inside each color gives the
matching signs in (P-F47a). This is exactly its exterior
coefficient contribution. Alternatively, first use
$A\to B$, then sort each column by color, and finally
transpose the grid. This gives
\[
 [e_{J_1}\wedge\cdots\wedge e_{J_d}]
                \theta_k^{\wedge n}
       =\varepsilon_T\tau_{n,k}\,I_T(F_n).
                                                        \tag{P-F47g}
\]
Both fixed signs are defined by specified permutations,
so (P-F47g) is an actual coefficient identity.
Sets of triples that are not $n$-regular have coefficient
zero, since a factorization into perfect matchings
would give degree $n$ at every vertex.

If the exterior power is nonzero, a nonzero coefficient
and (P-F47g) give the required invariant. Conversely,
the spanning statement for $I_T$ shows that a nonzero
invariant evaluation has a nonzero tableau evaluation.
The repeated-triple involution eliminates nonsimple
tableaux, and (P-F47g) then gives a nonzero exterior
coefficient. This proves (P-F47b) completely.

\paragraph{Exact nonvanishing ranges that follow immediately.}
For every $q$, $\theta_{3q}\ne0$. If $q$ is odd and
$n\ge2$, graded anticommutativity gives
$\theta_{3q}\wedge\theta_{3q}=0$, so its $n$th power
vanishes. This is consistent directly with Fermat:
a transposition of two variables fixes $F_n$ and has
determinant $-1$, whereas $k=3q$ is odd in (P-F47e).

For $q=2$, $k=6$, the twenty triples pair into ten
complementary pairs. Hence (P-F47a) has the form
\[
          \theta_6=\sum_{a=1}^{10}\epsilon_a u_a\wedge v_a,
             \qquad \epsilon_a\in\{1,-1\},              \tag{P-F47h}
\]
where the $u_a,v_a$ are a basis of $\bigwedge^3\mathbb C^6$.
Its $n$th power is nonzero exactly for $0\le n\le10$:
for $n\le10$ a product of any $n$ distinct pairs has
coefficient $n!\prod\epsilon_a\ne0$, and for $n>10$
the exterior degree exceeds twenty. Thus degree $2n$
invariants nonzero at Fermat exist exactly for $n\le10$.
This recovers the cubic case of B\"urgisser--Ikenmeyer,
Proposition~3.24, pp.~15--16; it is not a priority claim.

More generally the dimension bound gives
\[
      \theta_{3q}^{\wedge n}=0
           \quad\text{if}\quad
                  n>\binom{3q-1}{2}.                   \tag{P-F47i}
\]
For $n\le10$ and every even $q$, the nonzero degree-$2n$
invariant can be raised to the $q/2$ power; its Fermat
value stays nonzero. Equation (P-F47b) therefore proves
nonvanishing of $\theta_{3q}^{\wedge n}$ in this range.
These elementary arguments alone do not give general
nonvanishing for $n>10$ and even $q$.
P-F49 below supplies a polynomially bounded nonzero
power for every $n$, and P-F50 proves nonvanishing
for all sufficiently large even $q$, with a polynomial
threshold. The exact low-degree range remains distinct.

\paragraph{Remaining signed-factorization and index gates.}
Equation (P-F47g) asks for a simple regular triple
hypergraph whose perfect-matching factorizations have
nonzero signed sum. An unsigned existence theorem for
such factorizations cannot replace that signed assertion.
The solutions in P-F49--P-F50 give specified actual
nonzero coefficients and a polynomial conductor for
the Fermat degree monoid. They do not
by themselves control every residue-constituent polynomial
or the strict positivity and saturation indices in
GCT VI. Those additional implications require their
own proof. The primary locator for the determinant
constructions and their spanning property is
\url{https://arxiv.org/pdf/1511.02927}, equation~(3.7),
Theorem~3.11, and Remark~3.13, pp.~11--13.
